\documentclass[11pt]{article}
\usepackage[margin=2.3cm]{geometry}
\usepackage[T1]{fontenc}
\usepackage{lmodern}
\usepackage{longtable,booktabs,array}
\usepackage{listings}
\usepackage{xcolor}
\usepackage[colorlinks=true,linkcolor=blue!50!black,urlcolor=blue!50!black]{hyperref}
\lstset{basicstyle=\ttfamily\small,breaklines=true,columns=fullflexible,frame=single,
        backgroundcolor=\color{gray!6},rulecolor=\color{gray!40}}
\newcommand{\f}[1]{\texttt{#1}}
\setlength{\parindent}{0pt}\setlength{\parskip}{5pt}

\title{FOCEX agent: from a POSCAR to phonons\\[2pt]\large \f{FOCEX/AGENT/focex-agent.py} -- user manual}
\author{ALATDYN}
\date{October 2026}

\begin{document}
\maketitle
\tableofcontents

\section{What it does}
\f{FOCEX/AGENT/focex-agent.py} runs the FOCEX workflow in one of two ways:
\begin{itemize}
\item \textbf{full} (\f{-s POSCAR}): from the primitive cell of a crystal -- VASP relaxation and
  SCF, Born charges (DFPT), SC-SNAPS snapshots, the supercell VASP runs, the force collection,
  the five \f{*.params} files, FOCEX and the plots;
\item \textbf{FOCEX only} (\f{DATA: <folder>} in the instructions): from existing
  \f{POSCARi}/\f{FORCEDISPi} (and optionally Born charges) -- the \f{*.params} files, FOCEX and
  the plots (Section~\ref{sec:data}).
\end{itemize}
Every step uses defaults that can be overridden, and every step is checked and written to
\f{report.log}. The folder \f{FOCEX/AGENT} holds the script, this manual, an example input
(\f{focex-agent.example.txt}) and an end-to-end test without VASP (\f{tests/}).
\f{FOCEX/focex-agent.py} forwards to \f{AGENT/focex-agent.py}, so run scripts of older projects
keep working.

\section{Requirements}
\begin{itemize}
\item VASP-Flow (\url{https://github.com/KeivanS/VASP-Flow}) in \f{\textasciitilde/VASP-Flow}
  (or \f{\$VASP\_FLOW\_DIR}): the agent imports its INCAR/KPOINTS generator, GGA+U and MAGMOM
  rules, KPAR/NCORE rule, POTCAR builder, profiles and high-symmetry $k$-path. The POTCAR library
  is \f{\$VASP\_POTCAR\_DIR}; the VASP binary \f{\$VASP\_STD} (or the profile).
\item Python with numpy, spglib and pymatgen (VASP-Flow's environment, e.g.\
  \f{\textasciitilde/vaspenv/bin/python}).
\item \f{sc\_snaps.x} (SC-SNAPS), \f{read\_outcar.x} (\f{FOCEX/utility}), the FOCEX binary, and
  gnuplot for the plots.
\end{itemize}
The agent sets \f{OMP\_NUM\_THREADS=1} for the VASP it starts: an OpenMP-threaded BLAS inside
each MPI rank otherwise starts one thread per core in every rank (a 96-atom cell made no SCF step
in 30 minutes on 8 ranks).

\section{Running it}
\begin{lstlisting}
cp FOCEX/AGENT/focex-agent.example.txt instructions.txt   # Project and RANKS are required
~/vaspenv/bin/python FOCEX/AGENT/focex-agent.py init -i instructions.txt -s POSCAR --profile rivanna
cd <Project> && sbatch run_workflow.sh                   # bash run_workflow.sh on a workstation

# FOCEX only, on existing POSCAR1/FORCEDISP1: DATA: <folder> in instructions.txt
~/vaspenv/bin/python FOCEX/AGENT/focex-agent.py init -i instructions.txt
~/vaspenv/bin/python FOCEX/AGENT/focex-agent.py run <Project>

~/vaspenv/bin/python FOCEX/AGENT/focex-agent.py run <Project> --only snaps    # redo one task
~/vaspenv/bin/python FOCEX/AGENT/focex-agent.py run <Project> --retry-failed  # = resubmit_failed.sh
~/vaspenv/bin/python FOCEX/AGENT/focex-agent.py plots <Project>               # remake the plots
\end{lstlisting}
\f{run\_workflow.sh} calls \f{focex-agent.py run <Project>}, which does the tasks below in
order. A finished task is recorded in \f{state.json} and skipped, so after a time limit the same
script is simply submitted again.

\begin{longtable}{>{\ttfamily}l p{0.78\linewidth}}
\toprule
task & what it does\\ \midrule
relax & Niggli-reduced, oriented cell (\f{POSCAR.init}); VASP relaxation in \f{01\_relax},
  restarted from CONTCAR up to three times if EDIFFG is not reached\\
poscar0 & CONTCAR Niggli-reduced and oriented again: \f{POSCAR0}, the cell and Cartesian frame
  of every later step\\
scf & \f{02\_scf}: same INCAR (ISIF = 1, NSW = 0) and KPOINTS; reuses CHGCAR/WAVECAR when the
  basis is unchanged; band gap from EIGENVAL\\
dfpt & \f{02\_scf/dfpt}, only if there is a gap (\f{DFPT: auto}; \f{on}/\f{off}):
  $\varepsilon_\infty$ and Born charges\\
snaps & \f{03\_snaps}: \f{cell.inp}, \f{supercell.inp}, \f{snaps.inp}; SC-SNAPS; if the rms
  displacement is not within 15\,\% of \f{DISP\_TARGET} the temperature is rescaled and SC-SNAPS
  rerun (at most twice); the undistorted supercell is checked to be commensurate with POSCAR0\\
supercells & \f{04\_supercells/snap\_NNN}: INCAR as the SCF plus ISYM = 0, Gamma mesh
  (Section~\ref{sec:kmesh}), MAGMOM expanded to the supercell\\
vasp\_snaps & the supercell runs, one after the other; failures are diagnosed in \f{report.log}
  and \f{resubmit\_failed.sh} reruns them (an OUTCAR is never overwritten)\\
collect & \f{read\_outcar.x} in every folder: \f{FORCEDISP1}; \f{POSCAR1} = \f{poscar\_000};
  checks the net force and that the OUTCAR positions are those of the POSCAR\\
focex & \f{05\_focex}: the five \f{*.params} files and FOCEX; with Born charges also a
  \f{NOBORN} run; if equations/unknowns $<$ \f{LASSO\_RATIO} (2) also a \f{LASSO} run; the DOS
  range is reset to 1.1$\times$ the highest frequency when needed. A FOCEX run counts as failed
  if its exit code is not 0\\
plots & \f{bvel.plt}, \f{all.plt}, \f{qha.plt} with their placeholders filled and
  \f{makeplots.sh}, in every FOCEX folder\\
\bottomrule
\end{longtable}
For each FOCEX run the report gives the fit error, the numbers of equations and unknowns, the
lowest $\omega^2$ (a warning below $-25$\,cm$^{-2}$) and the highest frequency (a warning above
2500\,cm$^{-1}$).

\section{Defaults}
\begin{itemize}
\item \f{ACCURACY: Accurate} -- ENCUT 600\,eV, EDIFF $10^{-7}$, EDIFFG $-0.001$, PREC = Accurate
  in every step; \f{Normal} -- 450\,eV, $10^{-6}$, $-0.003$, PREC = Normal (cheaper, but the
  undistorted 96-atom ZnO supercell then kept forces of $\sim$0.01\,eV/\AA). DFPT uses
  EDIFF $10^{-8}$.
\item Relaxation and SCF $k$-mesh: VASP-Flow \f{coarse} (1000 $k$-points per reciprocal atom),
  \f{fine} (6000) if every element is a metal.
\item \f{SUPERCELL: SMALL} -- the multiple of the conventional cell with the fewest atoms whose
  smallest face distance is at least 10\,\AA, within 100 atoms (\f{LARGE}: 14\,\AA, 300 atoms).
\item 50 snapshots (\f{SNAPSHOTS}), 30\,K, 400\,cm$^{-1}$, target rms displacement 0.02\,\AA.
  The number of equations is $3N_{\rm atoms}$ per snapshot; 20 snapshots of 96 atoms give 5760
  equations, against 264 unknowns for ZnO with ranks 1--3.
\item \f{latdyn.params} mesh \f{fine}; \f{kpbs.params} from the automatic path.
\item \f{RANKS} has no default; \f{RANGE: 6 2 1 0 0 0 0} shells for ranks 2--8; \f{FC2RANGE: 0}.
\end{itemize}
\f{structure.params}: line 1 $a\,b\,c\,\alpha\,\beta\,\gamma$ of the conventional cell, line 2 the
primitive vectors in conventional units, line 3 \f{1.0 0}. POSCAR0 is rotated so that its
conventional cell lies exactly where FOCEX and SC-SNAPS build it ($a$ along $x$, $b$ in the $xy$
plane), so POSCAR1, FORCEDISP1 and the Born charges are all in that frame.

\section{Existing data (FOCEX only)}\label{sec:data}
Give \f{DATA: <folder>} in the instructions (no \f{-s}). The folder holds \f{POSCAR1},
\f{FORCEDISP1} (and \f{POSCAR2}, \f{FORCEDISP2}, \dots; all pairs are used) in any orientation,
and optionally
\begin{itemize}
\item \f{POSCAR}: the primitive cell in the frame of the data; otherwise spglib finds it from
  \f{POSCAR1}, keeping its orientation;
\item \f{dielectric.params}: $\varepsilon_\infty$ and the Born charges in the frame of the data,
  one block per atom of that \f{POSCAR} (without a \f{POSCAR} only with one atom per species).
\end{itemize}
\f{SPECIES: Mg O} names the species when the POSCARs have no element line. \f{init} symmetrizes
the lattice of the primitive cell (spglib's idealized lattice parameters, same basis and atom
order: lattice parameters given to 6--7 digits make FOCEX's symmetry analysis lose force-constant
groups), rotates it into the FOCEX frame, rotates \f{POSCARi} and the positions and forces of
\f{FORCEDISPi} with the same matrix and $\varepsilon$ and $Z$ as $R^{T}TR$, writes \f{POSCAR0}
and \f{04\_supercells/}, and marks relax \dots collect as provided; \f{run} then does
\f{focex} and \f{plots}. The report gives the rotation and the size of the symmetrization.

Checked on MgO data in a rhombohedral orientation: without Born charges the fit error (7.197\,\%),
the 35 unknowns, the highest frequency (647\,cm$^{-1}$), the bulk and shear moduli equal those of
a FOCEX run in the original frame.

\section{Supercell $k$-mesh}\label{sec:kmesh}
If the supercell matrix (rows, primitive units) is diagonal up to a permutation, the mesh is the
primitive one divided by the multiples, rounded up (6\,10\,8 with 2$\times$2$\times$2 gives
3\,5\,4). Otherwise -- every multiple of a centred conventional cell -- the supercell mesh keeps
the $k$-point spacing of the primitive mesh and is in general not exactly commensurate.
\f{KPOINTS\_snaps} sets it by hand.

\section{Project folder}
\begin{longtable}{>{\ttfamily}l p{0.62\linewidth}}
\toprule
instructions.txt & copy of the input (\f{init} appends the $k$-mesh density it chose)\\
POSCAR.input, POSCAR.init & the structure given; reduced and oriented\\
POSCAR0 & relaxed cell, frame of every later step\\
state.json & finished tasks, supercell matrix, gap, $\varepsilon$, Born charges\\
report.log & one entry per task: OK/FAILED, minutes, what was checked\\
run\_workflow.sh & the batch script (VASP-Flow profile); \f{resubmit\_failed.sh}\\
01\_relax, 02\_scf & VASP; \f{02\_scf/dfpt} the DFPT run\\
03\_snaps & SC-SNAPS input and \f{poscar\_NNN}\\
04\_supercells & \f{snap\_NNN/} one VASP run each; \f{FORCEDISP1}, \f{POSCAR1}\\
05\_focex & the FOCEX run and plots; \f{NOBORN/}, \f{LASSO/} when made\\
\bottomrule
\end{longtable}
WAVECAR, CHG, CHGCAR, WFULL, TMPCAR and PROCAR are deleted after use; OUTCARs never.

\section{Thermal conductivity}
THERMACOND reads \f{lat\_fc.dat}, \f{fc2.dat} and \f{fc3.dat} (and, if present,
\f{structure.params}: since October 2026 the cell is also taken from \f{lat\_fc.dat}); give at
least \f{RANKS: 1 1 1 0 0 0 0 0}. The agent does not run it: set lines 8--13 of
\f{latdyn.params} in \f{05\_focex} and run \f{kap8} there.

\section{Overriding any file}
A block named after a file replaces it:
\begin{lstlisting}
structure_params
... the whole file ...
structure_params_END
\end{lstlisting}
\f{INCAR} (every step) and \f{INCAR\_relax}, \f{INCAR\_scf}, \f{INCAR\_dfpt}, \f{INCAR\_snaps} are
merged tag by tag into the generated INCARs. The other blocks are written verbatim:
\f{KPOINTS}, \f{KPOINTS\_snaps}, \f{cell\_inp}, \f{supercell\_inp}, \f{snaps\_inp},
\f{structure\_params}, \f{latdyn\_params}, \f{kpbs\_params} (required with \f{KPATH: 1}),
\f{dielectric\_params}, \f{default\_params}. \f{END\_<NAME>} is also accepted, in any case.

\section{Born charges}
VASP prints $Z^*_{ij}$ with $i$ the field and $j$ the displacement direction. FOCEX contracts the
wave vector with the second index of \f{charge(i,j)} (\f{aux\_na.f90}), so the agent writes the
transpose; \f{BORN\_TRANSPOSE: no} writes them as printed.

\section{Testing without VASP}
\f{tests/run\_test.sh} runs the full workflow on fcc argon with \f{tests/fake\_vasp.py} (a
nearest-neighbour Lennard-Jones crystal) in place of VASP and compares FOCEX's phonons with the
exact ones (\f{tests/exact\_phonons.py}): within 0.23\,cm$^{-1}$ (0.4\,\%), elastic constants
within 0.8\,\%.

\end{document}
